\chapter{Reliability}

Reliability is about the system behaving well under stress.
It is a product feature, not just an operational concern.

\section{Timeouts}
Every network call must have a timeout.
No timeout means a slow failure becomes an outage.

\section{Retries}
Retries can help or hurt.
Only retry on safe operations and use exponential backoff.

\section{Circuit breakers}
Circuit breakers prevent repeated failures from cascading.
Open the circuit when a dependency becomes unreliable.

\section{Bulkheads}
Isolate critical resources so one component cannot starve another.
Use separate connection pools or worker queues.

\section{Graceful degradation}
When dependencies fail, degrade gracefully.
Serve cached data or reduced functionality.
Degradation keeps the core user path working.

\section{Error budgets}
Reliability goals should be measurable.
Define an error budget and track consumption.
If the budget is burned, slow releases and fix reliability gaps.

\section{Health checks}
Use liveness and readiness checks.
Readiness should verify critical dependencies.
Liveness should be minimal to avoid restart loops.

\section{Load shedding}
When overload hits, shed low-value traffic first.
Return fast errors for non-critical endpoints.
Load shedding prevents total collapse.

\section{Rate limiting}
Rate limits protect the system from spikes.
Use per-user and per-tenant limits when appropriate.
Expose limits in responses to support client backoff.

\section{Fallbacks}
Fallbacks are precomputed or cached responses.
Use fallbacks for non-critical data.
Document when a fallback is acceptable.

\section{Redundancy}
Single points of failure kill reliability.
Use redundant instances and multiple availability zones.
Redundancy should include databases and caches, not just stateless services.

\section{Disaster recovery}
Disaster recovery requires tested runbooks.
Backups are useless without restore drills.
Define RTO and RPO targets and measure them.

\section{Change management}
Most incidents come from changes.
Use canaries and progressive rollouts.
Rollback should be fast and rehearsed.

\section{Incident response}
Incident response needs clear roles.
Define incident commander, communications, and operations roles.
Practice incident response with game days.

\section{Post-incident reviews}
Post-incident reviews focus on learning, not blame.
Track corrective actions and verify completion.
Reliability improves only when fixes are implemented.

\section{Reliability metrics}
Track service availability, error rates, and latency.
Use SLOs to measure user impact.
Metrics should be tied to alerts with clear actions.

\section{Example reliability targets}
\begin{tabularx}{\textwidth}{@{}lXr@{}}
\toprule
Service & SLO & Target \\
\midrule
Task API & availability & 99.9\% \\
Billing API & availability & 99.95\% \\
Worker queue & job success rate & 99.5\% \\
\bottomrule
\end{tabularx}

\section{Chaos testing}
Chaos tests reveal hidden dependencies.
Run controlled failure scenarios in staging.
Use chaos tests to validate runbooks and alerts.

\section{Capacity management}
Reliability depends on headroom.
Keep enough capacity to handle spikes.
Capacity reviews should happen regularly.

\section{Example circuit breaker config}
\begin{lstlisting}[language=YAML]
circuit_breaker:
  failure_threshold: 0.2
  window_seconds: 60
  open_seconds: 30
\end{lstlisting}

\section{Timeout budgets}
Timeouts should align with latency budgets.
If a dependency times out after your SLO, it is too slow.
Timeouts should be short and explicit.

\section{Retry budgets}
Retries increase load and latency.
Define a retry budget so a spike does not overload dependencies.
Retry budgets are a key part of resilience engineering.

\section{Dependency isolation}
Critical dependencies should be isolated with bulkheads.
Non-critical dependencies should fail open when possible.
Isolation prevents cascading failures.

\section{Multi-region failover}
Failover should be automatic and tested.
DNS or traffic management must support fast switching.
Regular failover drills reduce surprises.

\section{Data durability}
Durability is part of reliability.
Use backups with verified restore procedures.
Measure backup success and restore time.

\section{Service ownership}
Reliability requires clear ownership.
Every service should have an on-call owner.
Ownership clarifies who responds during incidents.

\section{Monitoring maturity}
Start with core signals: latency, errors, saturation.
Add traces once logs and metrics are stable.
Instrumentation gaps should be tracked like bugs.

\section{Runbook example}
\begin{lstlisting}[language=YAML]
runbook:
  alert: "HighErrorRate"
  steps:
    - "check recent deploys"
    - "inspect dependency latency"
    - "rollback if needed"
\end{lstlisting}

\section{Failure mode catalog}
Maintain a catalog of known failure modes.
Document mitigations and detection signals.
This reduces time to diagnose repeat incidents.

\section{Reliability reviews}
Review reliability after major changes.
Confirm SLO impact and capacity headroom.
Reliability reviews prevent regressions.

\section{Dependency SLOs}
Dependencies need their own SLOs.
Track dependency error rates and latency.
If a dependency breaks its SLO, adjust retry and fallback behavior.

\section{Graceful degradation levels}
Define degradation levels to simplify incident response.
Each level should describe what features are disabled.
This prevents ad-hoc decisions under pressure.

\section{Degradation table}
\begin{tabularx}{\textwidth}{@{}lX@{}}
\toprule
Level & Behavior \\
\midrule
Level 0 & full functionality \\
Level 1 & disable non-critical features \\
Level 2 & serve cached data only \\
\bottomrule
\end{tabularx}

\section{Saturation signals}
Saturation includes CPU, memory, and queue depth.
Saturation signals are the earliest warning signs.
Alert on sustained saturation, not short spikes.

\section{Progressive delivery}
Rollouts should be progressive and reversible.
Use feature flags and staged rollouts.
Progressive delivery reduces incident blast radius.

\section{Synthetic monitoring}
Synthetic probes detect failures before users do.
Run key user flows on a schedule.
Treat synthetic alerts as early warnings.

\section{User-facing status}
Status pages build trust during incidents.
Update status pages with honest timelines.
User communication is part of reliability.

\section{Example rollout policy}
\begin{lstlisting}[language=YAML]
rollout:
  canary_percent: 5
  soak_minutes: 30
  expand_percent: 25
\end{lstlisting}

\section{Hedged requests}
Hedged requests send a second request when the first is slow.
They reduce tail latency at the cost of extra load.
Use them only for idempotent requests.

\section{Queue resilience}
Queues can buffer spikes but also hide issues.
Monitor queue age and failure rate.
If queues grow without bound, you are falling behind.

\section{Feature flags}
Feature flags enable rapid rollback.
Keep flags short-lived and documented.
Stale flags increase complexity and risk.

\section{Operational readiness}
Before launch, confirm runbooks and alerts exist.
Run a game day for major features.
Operational readiness reduces incidents after launch.

\section{Resilience scorecard}
Use a scorecard to track resilience practices.
Scorecards highlight gaps and drive improvements.

\section{Example resilience scorecard}
\begin{tabularx}{\textwidth}{@{}lX@{}}
\toprule
Practice & Status \\
\midrule
Timeouts on all deps & implemented \\
Retries with jitter & implemented \\
Runbooks per alert & partial \\
\bottomrule
\end{tabularx}

\section{Alert quality}
Alert quality is more important than quantity.
If an alert does not drive action, remove it.
Review alerts after each incident.

\section{Client timeouts}
Client timeouts must be shorter than server timeouts.
If clients wait longer, they retry and amplify load.
Set timeouts explicitly in SDKs.

\section{Dependency mapping}
Map critical dependencies for each service.
Dependency maps inform incident response.
Keep maps updated as systems evolve.

\section{Error budget policy}
Define actions when error budgets are spent.
Examples include feature freezes and reliability sprints.
Policies reduce debate during incidents.

\section{Example error budget policy}
\begin{lstlisting}[language=YAML]
policy:
  budget_burned: "pause feature releases"
  owner: "reliability"
\end{lstlisting}

\section{Incident communications}
Clear communication reduces confusion.
Provide incident updates on a schedule.
Use a template so updates stay consistent.

\section{Example incident update}
\begin{lstlisting}[language=YAML]
update:
  time: "10:30 UTC"
  impact: "increased latency"
  mitigation: "rolled back deploy"
\end{lstlisting}

\section{Failure modes}
Document common failure modes and mitigations.
If a failure repeats, add it to the catalog.
This shortens future recovery times.

\section{Failure mode table}
\begin{tabularx}{\textwidth}{@{}lX@{}}
\toprule
Failure & Mitigation \\
\midrule
DB saturation & shed load, scale replicas \\
Cache outage & bypass cache, reduce load \\
Queue backlog & scale workers, slow producers \\
\bottomrule
\end{tabularx}

\section{Capacity headroom}
Maintain headroom to handle spikes.
Headroom targets should be documented and monitored.
Running at 100\% capacity is a reliability risk.

\section{Autoscaling policies}
Autoscaling should be based on meaningful signals.
CPU alone is insufficient for queue-heavy systems.
Use queue depth and latency metrics for autoscaling.

\section{Traffic shaping}
Traffic shaping can smooth bursts.
Use token buckets for high-traffic endpoints.
Shaping reduces sudden load on dependencies.

\section{Postmortem template}
\begin{lstlisting}[language=YAML]
postmortem:
  summary: "what happened"
  impact: "user impact"
  root_cause: "primary cause"
  actions:
    - "short-term fix"
    - "long-term fix"
\end{lstlisting}

\section{SLO dashboards}
SLO dashboards should show current error budget burn.
Dashboards make reliability status visible to everyone.
Include links to runbooks and recent incidents.

\section{Dependency budgets}
Each dependency should have a latency and error budget.
Budgets help teams decide when to fail open or fail closed.
If a dependency exceeds its budget, degrade gracefully.

\section{Queue drain procedures}
Queues should be drained during maintenance.
Drain procedures prevent job loss during deployments.
Document drain steps in runbooks.

\section{Resilient client behavior}
Clients should retry safely with backoff.
Expose retry guidance in error responses.
Good client behavior reduces load during incidents.

\section{Latency budgets for dependencies}
Each dependency should have a latency budget.
Budgets enforce timeouts and fallback behavior.
Without budgets, slow dependencies silently erode reliability.

\section{Reliability guardrails}
Guardrails prevent risky changes from shipping.
Examples include SLO checks and mandatory canaries.
Guardrails reduce accidental regressions.

\section{Backpressure}
Backpressure keeps producers from overwhelming consumers.
Use bounded queues and explicit rejection when capacity is full.
If you cannot reject, you will eventually fail elsewhere.

\section{Cold start risk}
Cold starts increase latency and error rates.
Pre-warm caches and keep a minimum instance baseline.
Measure cold start impact in performance tests.

\section{Maintenance windows}
Plan maintenance windows for disruptive work.
Communicate clearly and use them sparingly.
Maintenance windows reduce surprise outages.

\section{Reliability debt}
Reliability debt accumulates when shortcuts are taken.
Track debt like bugs with owners and deadlines.
Ignoring debt increases the probability of major incidents.

\section{User impact mapping}
Map critical user journeys to system dependencies.
Use this map to prioritize alerts and SLOs.
Impact mapping keeps focus on what users feel.

\section{Stateful service resilience}
Stateful systems need extra care during failover.
Practice leader election and data rehydration in drills.
Failover plans should cover data integrity, not just uptime.

\section{Escalation policy}
Define clear escalation paths for incidents.
Escalate based on impact and time without mitigation.
Policies reduce confusion during high-severity events.

\section{Detection latency}
Reliability is limited by how fast you detect issues.
Measure time to detect and time to mitigate.
Short detection latency reduces user impact.

\section{Failure injection}
Failure injection tests real-world behavior.
Inject timeouts, errors, and latency into dependencies.
Use these tests to validate fallbacks and alerts.

\section{Reliability and cost}
Over-provisioning buys reliability but costs money.
Use error budgets to balance cost and risk.
The right balance depends on user expectations.

\section{On-call handoff}
On-call handoff should include current incidents and risks.
Handoffs reduce repeated work and missed context.
Document the handoff in a shared log.

\begin{checklistbox}{Checklist: reliability}
\begin{itemize}[nosep]
  \item Time out all dependency calls.
  \item Retry only idempotent operations.
  \item Use circuit breakers for flaky dependencies.
\end{itemize}
\end{checklistbox}
